\documentclass[11pt]{article}
\usepackage[margin=2.5cm]{geometry}
\usepackage{amsmath, amssymb}
\usepackage{booktabs}
\usepackage{hyperref}
\usepackage{parskip}

\title{Chapter 4 Architecture Decision Note\\
       \large Why the Cascade Controller Was Replaced by a Single-State Integrator}
\author{LLM Serving Control Project}
\date{\today}

\begin{document}
\maketitle

\section{Background}

Chapters 1--3 of this project applied classical cascade control to an LLM
inference serving system.  The inner loop regulated queue depth~$q$ by
adjusting batch size~$B$, and the outer loop regulated p95 latency~$\ell$
by setting the queue setpoint~$q_{\mathrm{ref}}$.  The plant model assumed:

\begin{equation}
    \ell[k] = \alpha B[k] + \gamma B[k]^2 + \beta\, q[k]
    \label{eq:model}
\end{equation}

with parameters $\alpha = 0.1$, $\gamma = 0.8$, $\beta = 2.0$ taken from a
theoretical stochastic simulation.  The key structural assumption was that
a persistent request queue~$q$ builds up and directly contributes to latency
via the term~$\beta q$.

\section{What We Found in Practice}

When the cascade controller was connected to a real \texttt{qwen2.5:3b}
model served by Ollama (Chapter~3), plant identification revealed two
critical discrepancies.

\subsection{The queue term is unidentifiable}

The Stage~2 identification experiment ran sustained Poisson arrivals at
$\lambda \in \{2, 3, 5, 7, 9\}$ requests per tick, with a drain rule
$B[k] = \min(B_{\max},\, q[k] + a[k])$.  In every case the measured
queue depth settled to $q \approx 0$ regardless of $\lambda$.

The root cause is how Ollama processes requests internally.  Unlike a
production serving engine such as vLLM---which maintains an explicit
scheduler queue and can hold requests pending while others are
processed---Ollama dispatches each HTTP request immediately to an available
inference slot.  Requests that arrive when all slots are occupied are queued
inside Ollama's internal HTTP stack, not in a user-visible counter.  From
MATLAB's perspective, the drain rule always matched arrivals, so the
software queue counter never accumulated.  The $\beta q$ term in
equation~\eqref{eq:model} has no observable counterpart in the real system.

\subsection{TTFT is a direct function of concurrency, not queue depth}

The Stage~1 identification swept concurrent batch size $B \in \{1, \ldots, 12\}$
by firing $B$ requests \emph{simultaneously} via \texttt{parfeval}.  The
measured Time to First Token (TTFT) followed a clear nonlinear curve:

\begin{equation}
    \ell(B) = c_0 + c_1 B + c_2 B^2
    \label{eq:fit}
\end{equation}

with fitted parameters:
\begin{align*}
    c_0 &= 97.16 \;\text{ms} \qquad \text{(idle prefill offset)} \\
    c_1 &= 30.35 \;\text{ms/req} \qquad \text{(linear contention)} \\
    c_2 &= -1.44 \;\text{ms/req}^2 \qquad \text{(saturation effect)}
\end{align*}

with $R^2 = 0.815$.  The negative $c_2$ reflects the fact that at high
concurrency Ollama spreads load across multiple GPU inference slots
(configured via \texttt{OLLAMA\_NUM\_PARALLEL=16}), partially counteracting
the linear contention.

The physical mechanism is GPU KV-cache memory bandwidth competition: when
$B$ sequences are processed in parallel, each competes for bandwidth, raising
TTFT.  This is a concurrency effect, not a queueing effect.

\subsection{The assumed $\beta$ was wrong by 8$\times$}

Linearising equation~\eqref{eq:fit} at the operating point $B_0 = 5$:

\begin{equation}
    \beta_{\mathrm{eff}} = \left.\frac{d\ell}{dB}\right|_{B_0}
    = c_1 + 2c_2 B_0 = 30.35 - 2(1.44)(5) = 15.97 \;\text{ms/req}
\end{equation}

The assumed $\beta = 2.0\;\text{ms/req}$ was $8\times$ too small.  The
outer-loop integral gain $K_{il} = (e^{-\Delta t/\tau} - 1)/\beta$ was
therefore $8\times$ too aggressive, causing the outer integrator to
over-correct on every tick.

\section{Revised Architecture: Single-State Integral Controller}

Since $q \approx 0$ always and the plant reduces to the static map
$\ell \approx f(B)$, the cascade architecture adds no value.  The inner
loop was regulating a state that does not exist in the real system.

The revised Chapter~4 controller is a single-state discrete-time integrator:

\begin{align}
    e[k]      &= \ell_{\mathrm{target}} - \ell_{\mathrm{meas}}[k] \\
    \xi[k+1]  &= \xi[k] + e[k] \label{eq:int}\\
    B[k]      &= \mathrm{clamp}\!\left(B_0 - K_{il}\,\xi[k],\; B_{\min},\; B_{\max}\right)
\end{align}

The closed-loop characteristic equation is:

\begin{equation}
    z_{\mathrm{cl}} = 1 + \beta_{\mathrm{eff}}\, K_{il}
\end{equation}

Setting $z_{\mathrm{cl}} = e^{-\Delta t / \tau_{\mathrm{cl}}}$ with
$\tau_{\mathrm{cl}} = 20\;\text{s}$ and $\Delta t = 1\;\text{s}$:

\begin{equation}
    K_{il} = \frac{e^{-\Delta t/\tau_{\mathrm{cl}}} - 1}{\beta_{\mathrm{eff}}}
           = \frac{e^{-0.05} - 1}{15.97}
           = -0.003054 \;\text{req/ms}
\end{equation}

The anti-windup limits on $\xi$ are derived from the actuator bounds:

\begin{equation}
    \xi \in \left[\frac{B_0 - B_{\max}}{K_{il}},\;
                   \frac{B_0 - B_{\min}}{K_{il}}\right]
         = [-1310,\; 2292]
\end{equation}

\section{Summary}

\begin{table}[h]
\centering
\begin{tabular}{lll}
\toprule
& \textbf{Chapter 3 (cascade)} & \textbf{Chapter 4 (single-loop)} \\
\midrule
Plant model       & $\ell = \alpha B + \gamma B^2 + \beta q$ & $\ell = f(B)$ (identified) \\
Controller states & $\xi_q$ (inner) + $\xi_l$ (outer)        & $\xi$ (one integrator) \\
Actuator          & $q_{\mathrm{ref}} \to B$ via inner loop   & $B$ directly \\
$\beta$ used      & 2.0 (assumed)                             & 15.97 (measured) \\
$K_{il}$          & $-0.0164$ ($8\times$ too aggressive)      & $-0.003054$ (correct) \\
Queue state       & Tracked in software                       & Not needed ($q \approx 0$) \\
\bottomrule
\end{tabular}
\caption{Architecture comparison between Chapter 3 and Chapter 4.}
\end{table}

The key insight is that \textbf{Ollama does not expose true request queueing
at the application level}.  TTFT is determined by GPU concurrency, not by
a waiting queue.  A single integral controller acting directly on batch size
is both sufficient and correct for this plant.  The cascade was solving a
problem that did not exist in the real system.

\end{document}
